\section{Diagonal covers and the complete scheme fibre}
\label{sec:diagonal-v171}
The diagonal calculation provides a second application and a comparison
with the tangent case. It also isolates the local algebra needed to
pass from generic multiplicities to a complete scheme fibre. Let
$b=u^\rho$, $c=v^\sigma$, where $\rho,\sigma$ are positive integers,
and let $Y=T_a\times_{S_a}\Spec k[u,v]$, $Z=Y^\nu$. Put
\begin{equation}\label{eq:diagonal-data-v171}
 g_i=\gcd(\rho,i\sigma),\quad
 n_i=(A_i,B_i)=\left(\frac{i\sigma}{g_i},\frac{\rho}{g_i}\right),
 \quad e_i=\frac{\rho\sigma}{g_i},\quad \ell_i=\min(A_i,B_i).
\end{equation}
Set $n_0=(0,1)$ and $n_{a+1}=(1,0)$, in clockwise order.

\begin{theorem}[Diagonal parameter boundary]
\label{thm:diagonal-v171}
The surface $Y$ is integral and a local complete intersection. Its
normalization is the normal toric surface with consecutive cones
$\langle n_i,n_{i+1}\rangle$. The scheme fibre $F$ of
$Z\to\Spec k[u,v]$ over the origin is Cohen--Macaulay of pure
dimension one, without embedded points. Its reduction is a nodal
chain of $a$ projective lines $D_i$. The generic length on $D_i$ is
$\ell_i$, and the nilradical has exact order $\max_i\ell_i$ (order
one means it is zero). The conductor is
\begin{equation}\label{eq:diagonal-conductor-v171}
 \mathfrak c_{Z/Y}=\OO_Z\left(-\sum_{i=1}^a\kappa_iD_i\right),
 \qquad
 \kappa_i=\frac{\rho i\sigma-\rho-i\sigma}{g_i}+1.
\end{equation}
There is one divisor over each $E_i$, with ramification index $e_i$
and residue degree $g_i$. In particular $Y$ is normal if and only if
$\rho=1$.
\end{theorem}

\subsection{The raw charts and the graded semigroup}
Flat base change of the exhaustive cover of $T_a$ gives the rings
\begin{align}
 A_0&=k[u,v,k_0]/(v^\sigma-u^\rho k_0),\label{eq:raw-first-v171}\\
 A_i&=k[u,v,e,h]/(u^\rho-e^{i+1}h^i,\ v^\sigma-eh)
       \quad(1\leq i<a),\label{eq:raw-middle-v171}\\
 A_a&=k[u,v,w]/(u^\rho-v^{a\sigma}w).
 \label{eq:raw-last-v171}
\end{align}
The middle ring is finite free over $k[e,h]$, with basis
$u^jv^l$, $0\leq j<\rho$, $0\leq l<\sigma$. The end rings are
similarly the pullbacks of the end charts. They inject into the
single generic field $k(u,v)$ by finite flatness over the respective
coordinate domains of $T_a$. Thus these equations have neither a
hidden vertical component nor an embedded component. They cover $Y$,
not just a surface birational to $Y$.

For the exponents $E_j$ of Proposition~\ref{prop:chain-factorization-v171}
put
\[
 I=\prod_{j=1}^a(u^\rho,v^{j\sigma})^{E_j},\qquad
 \eta_i=\sum_{j=1}^aE_j\min(i,j).
\]
Let $H\subset\mathbb Z^3$ be the semigroup generated by $(1,0,0)$,
$(0,1,0)$ and $(\alpha,\beta,1)$ for the monomial generators
$u^\alpha v^\beta$ of $I$. Its algebra is the ordinary Rees algebra.
The group generated by $H$ is all $\mathbb Z^3$: subtracting the
first two coordinate vectors from any degree-one generator gives
$(0,0,1)$ in the group. The rational cone of $H$ has the complete
half-space description
\begin{equation}\label{eq:semigroup-facets-v171}
 \alpha,\beta,n\geq0,\qquad
 i\sigma\alpha+\rho\beta\geq n\rho\sigma\eta_i
       \quad(1\leq i\leq a).
\end{equation}
Indeed its degree-one section is the Minkowski sum of the Newton
segments of the factors plus the nonnegative quadrant; each compact
edge has inward normal $(i\sigma,\rho)$ and positive length.
The coordinate and degree facets account for its noncompact faces.
The normalization is therefore
\begin{equation}\label{eq:normal-rees-all-degrees-v171}
 \overline{\mathcal R(I)}
 =k[u^\alpha v^\beta Q^n:(\alpha,\beta,n)\in
                       \mathbb Z^3\text{ satisfying }
                       \eqref{eq:semigroup-facets-v171}].
\end{equation}
To justify the grading in this assertion, a lattice point in the cone
is a rational nonnegative combination of the listed generators.
Multiplying by a common denominator puts a positive multiple in
$H$, so its monomial is integral, satisfying a monic equation
$T^N-u^{N\alpha}v^{N\beta}Q^{Nn}$. Conversely an integral monomial
lies in the cone by every facet valuation. Thus the saturation is
$\operatorname{Cone}(H)\cap\mathbb Z^3$ with its unchanged third-coordinate
grading, not a claim that the original semigroup was saturated in a
fixed degree. This is the standard affine semigroup-normalization
theorem; see \cite[Chapters 1 and 5]{HunekeSwanson171}.

The raw evaluation ideal is $u^{\rho C}I$. The degree-preserving
isomorphism $fQ^n\mapsto u^{n\rho C}fQ^n$ identifies its Rees algebra
with $\mathcal R(I)$ and, on their fraction fields, identifies the
normalizations by $Q\mapsto u^{\rho C}Q$. This distinguishes a graded
presentation from its Proj. Formula \eqref{eq:normal-rees-all-degrees-v171}
gives the fan in Theorem~\ref{thm:diagonal-v171}.

\subsection{A coordinate ideal on an arbitrary proper cone}
\begin{lemma}[The full coordinate ideal]
\label{lem:coordinate-ideal-v171}
Let $l,h$ be clockwise primitive rays spanning a two-dimensional
cone strictly smaller than the nonnegative quadrant, and lying in
that quadrant. In the character lattice $M=\mathbb Z^2$ put
\[
 A=k[\{m\in M:\langle l,m\rangle,\langle h,m\rangle\geq0\}].
\]
The characters $u=\chi^{(1,0)}$, $v=\chi^{(0,1)}$ belong to $A$, and
\begin{equation}\label{eq:coordinate-divisorial-v171}
 (u,v)A=\bigoplus_{\substack{m\in M\,;\,
 \langle l,m\rangle\geq\min(l_1,l_2)\\
 \langle h,m\rangle\geq\min(h_1,h_2)}} k\chi^m.
\end{equation}
It is the intersection of the two height-one symbolic valuation
ideals with these orders. In particular it is reflexive, and
$A/(u,v)A$ has no zero-dimensional associated prime whenever it is
nonzero.
\end{lemma}
\begin{proof}
Negative exponents are allowed exactly when both dual-cone
inequalities hold. They are characters regular on this chart,
although they need not be polynomials in the original plane. Divisibility
by $u$ means $m-(1,0)$ satisfies both inequalities; divisibility by
$v$ means $m-(0,1)$ does. The forward inclusion in
\eqref{eq:coordinate-divisorial-v171} follows immediately.

If $l_1\leq l_2$ and $h_1\leq h_2$, the character $v/u$ is regular,
so $(u,v)=(u)$ and the reverse inclusion follows. If both inequalities
are reversed, use $(u,v)=(v)$. Equality of one pair of coordinates
causes no difficulty in these cases.

In the remaining case $l_1<l_2$ and $h_1>h_2$. Write $m=(x,y)$,
$L=x+y$. We first prove $L\geq1$ for an integral $m$ on the right
of \eqref{eq:coordinate-divisorial-v171}. If $L\leq0$, the low-ray
inequality
$l_1L+(l_2-l_1)y\geq l_1$ forces $y\geq0$, with $y>0$ if $l_1>0$;
the high-ray inequality
$h_1L-(h_1-h_2)y\geq h_2$ forces $y\leq0$, with $y<0$ if $h_2>0$.
At least one of $l_1,h_2$ is positive, because the full quadrant is
excluded. This is a contradiction.

Subtraction of $(1,0)$ can fail only at the high ray: the low
threshold is already $l_1$. Subtraction of $(0,1)$ can fail only
at the low ray: the high threshold is already $h_2$. If both fail,
then
\[
 l_1L+(l_2-l_1)y<l_2,
 \qquad h_1L-(h_1-h_2)y<h_1.
\]
Since $L\geq1$ and $y$ is integral, the first inequality gives
$y\leq0$, while the second gives $y>0$. This contradiction proves
that one of $u,v$ divides $\chi^m$ in $A$.

In a normal affine toric ring the two ray valuations are the
height-one valuations of its invariant prime divisors. The monomial
ideal on the right is their symbolic-ideal intersection; orders at
all other height-one primes are nonnegative automatically, since
each character is a unit on the coefficient torus. Consequently it
is the divisorial ideal specified by these orders and is reflexive.
The normal surface ring is $S_2$, and a reflexive ideal has depth
two at each closed point. In
$0\to (u,v)A\to A\to A/(u,v)A\to0$, the depth lemma therefore gives
depth at least one for the quotient at every closed point in its
support. All nonempty support has dimension one, so there is no
zero-dimensional associated prime.
\end{proof}

\begin{proof}[Proof of the fibre assertions in Theorem~\ref{thm:diagonal-v171}]
Every maximal cone in the fan is proper because $a\geq2$. Apply
Lemma~\ref{lem:coordinate-ideal-v171} to obtain, on the whole surface,
\[
 (u,v)\OO_Z=\OO_Z\left(-\sum_i\ell_iD_i\right).
\]
This is an equality of actual ideals. It proves purity,
Cohen--Macaulayness and the generic lengths by localization at the
DVR of $D_i$. The radical is the divisorial ideal with order one
on each $D_i$. If $L=\max\ell_i$, every ordinary product of $L$
sections of the radical has order at least $L$ on all $D_i$ and
hence belongs to $(u,v)\OO_Z$. This proves the upper bound on the
ordinary power of the nilradical, without asserting that an ordinary
product of reflexive ideals is reflexive. At a component with
$\ell_i=L>1$, the quotient of a DVR by its $L$th uniformizer power
has nonzero $(L-1)$st nilradical power, proving exactness.

At a crossing choose the standard finite cyclic cover of the toric
chart. If its adjacent determinant is $\Delta$, its ring is
$k[x,y]^{\mu_\Delta}$, with weights coprime to $\Delta$ in each
coordinate. The reduced boundary union is the invariant ring of
$k[x,y]/(xy)$, since invariants are exact in characteristic zero.
The latter invariant ring is
\[
 k[x^\Delta,y^\Delta]/(x^\Delta y^\Delta).
\]
Equivalently it is $k[x^\Delta]\times_k k[y^\Delta]$. Thus the
completed node ring is $k[[X,Y]]/(XY)$ even if the ambient toric
surface is singular. At end charts only one boundary axis is in the
fibre. The exceptional orbit closures are projective lines by the
complete one-dimensional star fans, so they glue to the claimed chain.
\end{proof}

\subsection{The conductor and its extension across intersections}
The diagonal map has Jacobian ideal
$(u^{\rho-1}v^{\sigma-1})$. Over $E_i$, the primitive normalized
valuation is $n_i$; its ramification index is $e_i$ in
\eqref{eq:diagonal-data-v171}. The residue character
$\zeta=u^{\rho/g_i}/v^{i\sigma/g_i}$ satisfies
$\zeta^{g_i}=b/c^i$. This polynomial is irreducible over
$k(b/c^i)$, so there is one divisor with residue degree $g_i$;
$e_i g_i=\rho\sigma$. Theorem~\ref{thm:duality-conductor-v171}
now gives
\[
 (\rho-1)A_i+(\sigma-1)B_i-e_i+1
 =\frac{\rho i\sigma-\rho-i\sigma}{g_i}+1,
\]
including every intersection as an equality of conductor ideals.
For reference, the following general extension argument also explains
why a separate hidden punctual condition is impossible here.

\begin{lemma}[$S_2$ extension of conductor conditions]
\label{lem:conductor-depth-v171}
Let $A$ be an $S_2$ Noetherian domain with finite normalization $B$.
Then $A=\bigcap_{\operatorname{ht}P=1}A_P$ in its fraction field, and
\begin{equation}\label{eq:conductor-intersection-v171}
 (A:B)=B\cap\bigcap_{\operatorname{ht}P=1}
             (A_P:B\otimes_A A_P).
\end{equation}
When the local conductor is written by valuations in the semilocal
normalization of $A_P$, the intersection includes each of its maximal
ideals, namely the primes of $B$ above $P$.
\end{lemma}
\begin{proof}
For the first assertion take a fraction $z$ in all $A_P$ and consider
the finite module $M=(A+Az)/A$. If it is nonzero, a minimal prime of
its support has height at least two. Localize there; $M$ then has
finite length, $A$ has depth at least two, and the torsion-free
module $A+Az$ has depth at least one. The exact sequence
\[
 0\longrightarrow A\longrightarrow A+Az\longrightarrow M
 \longrightarrow0
\]
and the depth inequality
$\operatorname{depth}M\geq\min(\operatorname{depth}A-1,
\operatorname{depth}(A+Az))$ give depth at least one, contrary to
nonzero finite length. Thus $M=0$. Testing $zb\in A$ for each $b\in B$
gives \eqref{eq:conductor-intersection-v171}.

Localization of finite normalization gives a semilocal normal
one-dimensional ring over $A_P$. Its maximal ideals are not to be
identified with $P$ itself; all primes above $P$ must be used.
Every nonzero conductor ideal of this semilocal Dedekind ring is
specified by its orders in these DVRs. Intersecting their conditions
inside $B$ proves the assertion, with no additional condition in
codimension two. In the geometric applications all rings are excellent,
so computing these orders after faithfully flat completion is legitimate;
in the present paper they were computed directly by trace and valuation.
\end{proof}

The coefficients in \eqref{eq:diagonal-conductor-v171} vanish when
$\rho=1$. If $\rho>1$, the last is positive, since $a\sigma\geq2$:
for $R,Q\geq2$, $RQ-R-Q+\gcd(R,Q)>0$. This proves the normality
criterion. Its asymmetry records the orientation of the chain
$\prod_j(b,c^j)$, not a symmetry-breaking assertion about arbitrary
covers of the plane.

If $\rho=\sigma=d$, the conductor generators on the first and every
intermediate normalized chart are $u^{d-1}$, and on the last chart
$v^{a(d-1)}$. Indeed their valuations are the conductor coefficients
on all rays of that chart, including order zero on its nonexceptional
axis; the ideal identity follows from the theorem. At an intermediate
intersection the same character is used on both sides. At the last
intersection $u=v^a\zeta$, where $\zeta$ is a unit, so the generators
are related by the unit $\zeta^{d-1}$. This verifies their gluing,
not just equality of divisorial orders on one chart.

\subsection{Quotient conventions and an unequal-cover chart}
To specify the quotient type without an orientation ambiguity, order
the two primitive rays $r_1,r_2$ counterclockwise and put
$\Delta=\det(r_1,r_2)>0$. Extend $r_1$ to an oriented lattice basis
$(r_1,t)$ and write $r_2=pr_1+\Delta t$. Then the chart has cyclic
type $\frac1\Delta(-p,1)$, where $p$ is read modulo $\Delta$ and
$\gcd(p,\Delta)=1$. This follows because
$t=(-p r_1+r_2)/\Delta$ acts on the smooth covering characters with
these weights. Interchanging the two rays interchanges the weights;
changing $t$ changes $p$ by a multiple of $\Delta$.

\begin{example}[An unequal cover at a quotient crossing]
\label{ex:unequal-cover-v171}
Take $(a,\rho,\sigma)=(3,2,3)$. On the first raw chart
\[
 A=k[u,v,\theta]/(v^3-u^2\theta)
\]
the normalization is the cubic Veronese ring
\[
 B=k[s^3,s^2t,st^2,t^3]
   =k[u,v,w,\theta]/(v^2-uw,\ vw-u\theta,\ w^2-v\theta).
\]
The map is $u=s^3$, $v=s^2t$, $w=st^2$, $\theta=t^3$.
The extra generator $w=v^2/u$ is integral, and the Veronese ring is
normal (it is the invariant ring for the scalar $\mu_3$ action).
It has the same fraction field and is finite, so is the normalization.
The cone rays are $(0,1),(3,2)$ and the quotient type is
$\frac13(1,1)$.

Here $B=A+Aw$, $w^2=v\theta$, and $uw=v^2$, $vw=u\theta$.
Thus $(u,v)B$ is contained in the conductor. Conversely the valuation
of $w$ on the exceptional ray is one, while $u,v,\theta$ have
values $3,2,0$. The conductor coefficient is two by
\eqref{eq:diagonal-conductor-v171}; every invariant monomial of
order at least two is divisible by $u$ or $v$, since $w^2=v\theta$.
Consequently
\[
 \mathfrak c_{B/A}=(u,v)B,\qquad
 B/(u,v)=k[w,\theta]/(w^2).
\]
This displays both the conductor and the nonreduced full fibre at the
quotient point. The fibre has generic length two and no embedded
associated prime. On the complete surface the three conductor orders
are $2,3,8$, while the three generic fibre lengths are $2,1,2$;
these are different invariants.
\end{example}

\subsection{Marked logarithmic compatibility}
For the diagonal family the normal fan is the coarsest subdivision
on which each $\min(\rho p,j\sigma q)$ is linear. Equivalently it
is terminal among normal integral tests with dense coefficient torus
which make every $(u^\rho,v^{j\sigma})$ invertible, with the category
specified as in Proposition~\ref{prop:chain-category-v171}.
For ordered transverse Cartier divisors $U,V$ on a smooth surface,
these ideals are intrinsically
$\OO(-\rho U)+\OO(-j\sigma V)$. Unit changes of local equations
preserve them; normalized blow-ups consequently glue. Subsequent
root covers compose by integral closure in the same function field.
The conductor-divisor transitivity is the special case of
Theorem~\ref{thm:conductor-tower-v171}, which also applies when the
coefficient change is not diagonal.

If $\partial Z$ is the full reduced toroidal boundary and
$\partial S'$ is the union of the coordinate axes, then
\[
 K_Z+\partial Z=\pi^*(K_{S'}+\partial S').
\]
Both sides are Cartier (indeed represented by the divisor of the
logarithmic torus form); $Z$ itself is simplicial, so its invariant
Weil divisors are $\mathbb Q$-Cartier. A toric resolution preserves
this equality and has a simple normal crossing reduced boundary,
proving log canonicity of the pair. This is a statement about the
parameter space, not stable or semistable reduction of its universal
embedded curves.
